\documentclass[11pt]{article}
\usepackage[margin=1in]{geometry}
\usepackage{amsmath,amssymb,amsthm,booktabs,array}
\usepackage[T1]{fontenc}
\usepackage{lmodern}
\usepackage[colorlinks=true,allcolors=blue,breaklinks=true]{hyperref}
\hypersetup{pdftitle={A converse obstruction and proofs of Paul D. Hanna's self-substitution conjectures after Li--Zhong--Zhou},pdfauthor={Haobo Ma and Wenlin Zhang}}
\newtheorem{theorem}{Theorem}
\newtheorem{lemma}[theorem]{Lemma}
\newtheorem{corollary}[theorem]{Corollary}
\newcommand{\Z}{\mathbb Z}
\newcommand{\Q}{\mathbb Q}
\newcommand{\ord}{\operatorname{ord}}
\title{A converse obstruction and proofs of Paul D. Hanna's self-substitution conjectures after Li--Zhong--Zhou}
\author{Haobo Ma\\\small ChronoAI Pte Ltd, Singapore
\and Wenlin Zhang\\\small National University of Singapore\\\small ChronoAI Pte Ltd, Singapore}
\date{17 November 2026}
\begin{document}
\maketitle
\begin{abstract}
We give alternative proofs of Paul D. Hanna's conditional integrality
conjecture for rational sequences with integral exponential input, in the
unpaginated comment of OEIS A163921, revision \#2
(30 March 2012), following Li--Zhong--Zhou's polynomial integrality theorem
(\nolinkurl{arXiv:2611.16283v1}, 12 November 2026, Theorem 1.2).
An infinite family of integral fixed points with nonintegral logarithmic
inputs shows that the converse, which Hanna did not assert, fails.
We give an explicit reduction of a signed weighted formulation to that
theorem, as well as a direct proof by Frobenius congruences and formal
iteration. The same methods establish the shifted equations of A156907
and A156908. An Euler product proves the sparse equation of A179471.
All series are ordinary formal power series.
\end{abstract}

\section{The question and a weighted theorem}

For $L:\Z_{>0}\to\Q$, Hanna's comment asserts the following ``Conjecture'':
\begin{equation}\label{input}
 F_L(x)=\exp\left(\sum_{n\geq1}\frac{L(n)}{n}x^n\right)\in\Z[[x]]
 \quad\Longrightarrow\quad
 G(x)=\exp\left(\sum_{n\geq1}\frac{x^n}{n}G(L(n)x)^n\right)\in\Z[[x]].
\end{equation}
This is the comment of A163921~\cite{oeis163921}; the entry itself takes
$L(n)=\sigma(n)=\sum_{d\mid n}d$. The related A157675 takes $L(n)=2^n$
and explicitly conjectures integrality~\cite{oeis157675}.
We give alternative proofs of both assertions and uniqueness, using and
comparing with Li--Zhong--Zhou~\cite[Theorem 1.2]{lzz}.
A weaker condition on $L$ suffices, and allowing weights gives two further
conjectures of Hanna. The converse obstruction below separates integral
fixed points from integral logarithmic inputs.

Write $[x^j]B$ for the coefficient of $x^j$ in a series $B$.
All infinite sums and products below are formal: in each degree only finitely
many terms contribute. A congruence of series means a congruence of every
coefficient. For a prime $p$ and a positive integer $m$, let $v_p(m)$ be
the exponent of $p$ in $m$.

\begin{theorem}\label{weighted}
Let $L,W:\Z_{>0}\to\Z$. Suppose that, for every prime $p$ and $m\geq1$,
\begin{equation}\label{congruences}
 L(pm)\equiv L(m)\pmod p,
 \qquad W(pm)\equiv W(m)\pmod {p^{v_p(m)+1}}.
\end{equation}
Then the equation
\begin{equation}\label{weighted-equation}
 G(x)=\exp\left(\sum_{n\geq1}\frac{W(n)}{n}x^nG(L(n)x)^n\right)
\end{equation}
has exactly one solution in $\Q[[x]]$, and it belongs to $1+x\Z[[x]]$.
Existence and uniqueness over $\Q$ hold for arbitrary rational $L,W$.
\end{theorem}

The tools are classical. The local exponential criterion is the
Dieudonn\'e--Dwork lemma; a precise formulation appears in
Krattenthaler--M\"uller~\cite[Lemma 1]{km}. Integer Euler products and their
divisor sums are the usual Euler-transform description of Dold sequences;
see Byszewski--Graff--Ward~\cite[Theorem 2.5 and Lemma 2.6]{bgw}.
Pomerat--Straub~\cite[Theorem 1.1 and Corollary 4.2]{ps} discuss the
neighboring problem of integral rational powers of a given series and
the same exponential criterion. Here the unknown series also occurs
inside its own dilated powers. We give the needed arguments below. The step specific to
\eqref{weighted-equation} is that the outer power $pm$ amplifies the
Frobenius congruence enough to cancel the denominator $m$.

\subsection{Relation to Li--Zhong--Zhou's theorem}

Li--Zhong--Zhou~\cite[Theorem 1.2, equations (1.5)--(1.6),
and Section 5]{lzz} prove that polynomials $R_n(y)\in\Z[y]$ satisfying
$R_{pm}(y)\equiv R_m(y)\pmod p$ give
\[
 \exp\left(\sum_{n\geq1}R_n(y)^n\frac{t^n}{n}\right)
 \in1+t\Z[y][[t]].
\]
Their stated fixed-point equation substitutes a single series for $y$.
The dilations and weights in Theorem~\ref{weighted} follow by the following
Euler-product reduction, rather than by identifying those two equations.

Put $c_d=d^{-1}\sum_{e\mid d}\mu(d/e)W(e)$, where $\mu$ is the
M\"obius function. If $d=p^rh$ with $p\nmid h$, then
\[
 dc_d=\sum_{e\mid h}\mu(h/e)
       \bigl(W(p^re)-W(p^{r-1}e)\bigr)\in p^r\Z.
\]
Thus $c_d\in\Z$ and $W(n)=\sum_{d\mid n}dc_d$.
For a polynomial $B\in1+y\Z[y]$, take
$R_m^{(d)}(y)=B(L(dm)y)^d$ in their theorem. These polynomials satisfy
its congruence hypothesis, so
\[
 H_d(t,y)=\exp\left(\sum_{m\geq1}
       B(L(dm)y)^{dm}\frac{t^m}{m}\right)\in1+t\Z[y][[t]].
\]
Taking logarithms and putting $n=dm$ gives
\[
 \prod_{d\geq1}H_d(x^d,x)^{c_d}
 =\exp\left(\sum_{n\geq1}\frac{W(n)}n x^nB(L(n)x)^n\right).
\]
The product is integral, including negative $c_d$, and is well-defined
degree by degree. Truncating $B$ extends the identity and integrality to
$B\in1+x\Z[[x]]$. The coefficient-stabilization argument below gives
existence and uniqueness of the fixed point. Hence the full signed weighted
theorem, and consequently Hanna's original assertion, follow from
Li--Zhong--Zhou's polynomial theorem. The local proof below gives a direct
alternative, with the same classical power-lifting mechanism.

For the shifted family one can also apply their theorem directly with
$R_n(y)=a^nB(a^ny)$: Fermat's congruence checks its hypothesis for every
integer $a$, and $R_n(y)^n=a^{n^2}B(a^ny)^n$.
The scalar weights $a^{n^2}$ already have an integrality proof for every
integer $a$ in Bala's A155200 comment (25 December 2019)~\cite{bala}.
For positive bases, Jaidee--Moss--Ward~\cite[Theorem 2 and Example 7]{jmw}
give the more general power-time-change result, crediting Moss's 2003 work.
These scalar results are ingredients, not assertions of the dilated
self-substitution equation.

\section{Local integrality and the proof}

Fix a prime $p$ and put
$R_p=\Z_{(p)}=\{a/b:a,b\in\Z,\ p\nmid b\}$.
The elementary identity
$\Q\cap\bigcap_p R_p=\Z$ is understood inside $\Q$.

\begin{lemma}[Power lifting]\label{lifting}
If $U,V\in R_p[[x]]$ and $U-V\in pR_p[[x]]$, then, for $m\geq1$,
\[
 U^m-V^m\in p^{v_p(m)+1}R_p[[x]].
\]
\end{lemma}
\begin{proof}
If $U-V=p^rD$ with $r\geq1$, expand $(V+p^rD)^p-V^p$.
The terms with exponents $1,\ldots,p-1$ contain a factor $p^{r+1}$;
the last contains $p^{rp}$, and $rp\geq r+1$, including $p=2$.
Write $m=p^sb$ with $p\nmid b$. Applying this observation $s$ times
shows that $U^{p^s}-V^{p^s}$ is divisible by $p^{s+1}$.
Raising these two series to the power $b$ preserves that divisibility,
by the factorization of a difference of $b$th powers.
\end{proof}

\begin{lemma}[Local exponential criterion]\label{dwork}
For $h\in x\Q[[x]]$, if
\[
 ph(x)-h(x^p)\in pxR_p[[x]],
\]
then $\exp(h)\in1+xR_p[[x]]$.
\end{lemma}
\begin{proof}
Put $E=\exp(h)$ and $D=ph(x)-h(x^p)$.
For $j\geq1$ one has $v_p(j!)\leq j-1$, since
$v_p(j!)=\sum_{r\geq1}\lfloor j/p^r\rfloor<j$.
Thus $\exp(D)\in1+pxR_p[[x]]$ and
\[
 E(x)^p/E(x^p)=\exp(D)\in1+pxR_p[[x]].
\]
Induct on the degree $N$ to prove that $e_N=[x^N]E$ lies in $R_p$.
Assume this below degree $N$ and set
$B=1+\sum_{1\leq j<N}e_jx^j$. Frobenius in the residue field $\mathbb F_p$
gives $B(x)^p\equiv B(x^p)\pmod p$. Both series are units of $R_p[[x]]$,
so $B(x)^p/B(x^p)\in1+pxR_p[[x]]$.
Consequently $Q=E/B$ satisfies $Q(x)^p/Q(x^p)\in1+pxR_p[[x]]$.
But $Q=1+e_Nx^N+O(x^{N+1})$, so the coefficient of $x^N$ in this
ratio is $pe_N$. It lies in $pR_p$, hence $e_N\in R_p$.
\end{proof}

\begin{proof}[Proof of Theorem~\ref{weighted}]
For $B\in1+x\Q[[x]]$, define
\[
 H_B(x)=\sum_{n\geq1}\frac{W(n)}n x^n B(L(n)x)^n,
 \qquad T(B)=\exp(H_B).
\]
First suppose $B\in1+xR_p[[x]]$ and \eqref{congruences} holds.
Pair the terms of index $pm$ in $pH_B(x)$ with those of index $m$
in $H_B(x^p)$ to obtain
\begin{align}\label{pairing}
 pH_B(x)-H_B(x^p)
 ={}&\sum_{p\nmid n}\frac{pW(n)}n x^nB(L(n)x)^n\notag\\
 &+\sum_{m\geq1}\frac{x^{pm}}m
 \bigl(W(pm)B(L(pm)x)^{pm}-W(m)B(L(m)x^p)^m\bigr).
\end{align}
The first sum belongs to $pxR_p[[x]]$.
For the second put $U=B(L(pm)x)^p$ and $V=B(L(m)x^p)$.
Reduction modulo $p$, Frobenius, and $L(pm)^p\equiv L(m)\pmod p$ give
\[
 U\equiv B(L(pm)^p x^p)\equiv B(L(m)x^p)=V\pmod p.
\]
The reduction of a coefficient of $B$ is fixed by its $p$th power,
because it lies in $\mathbb F_p$. If $s=v_p(m)$, Lemma~\ref{lifting}
therefore makes $U^m-V^m$ divisible by $p^{s+1}$.
The bracket in \eqref{pairing} equals
\[
 W(pm)(U^m-V^m)+(W(pm)-W(m))V^m,
\]
whose two terms are divisible by $p^{s+1}$.
Division by $m$ leaves a multiple of $p$. Lemma~\ref{dwork} now
shows that $T(B)\in1+xR_p[[x]]$.
Applying this at every prime shows that $T$ preserves $1+x\Z[[x]]$.

For existence and uniqueness, if $B-C\in x^d\Q[[x]]$ with $d\geq1$,
then $H_B-H_C\in x^{d+1}\Q[[x]]$: dilation and taking a positive
power preserve divisibility by $x^d$, and each summand has a further
factor $x^n$, $n\geq1$. Taking the exponential preserves that
divisibility as well. Hence
\[
 B\equiv C\pmod{x^d}\quad\Longrightarrow\quad
 T(B)\equiv T(C)\pmod{x^{d+1}}.
\]
Starting at $G_0=1$ and setting $G_{r+1}=T(G_r)$, it follows inductively
that $G_{r+1}\equiv G_r\pmod{x^{r+1}}$.
Each coefficient stabilizes, defining $G=\lim_rG_r$ coefficientwise.
The finite dependence of each coefficient gives $T(G)=G$.
Every solution has constant coefficient $1$ and agrees successively
with $G$ modulo $x,x^2,\ldots$, proving uniqueness.
Under \eqref{congruences}, every $G_r$ is integral, so $G$ is integral.
\end{proof}

\section{Hanna's conjectures and related entries}

\begin{corollary}\label{original}
The implication \eqref{input} holds, and its fixed point is unique.
In fact, integer $L$ and the first congruence of \eqref{congruences}
already suffice.
\end{corollary}
\begin{proof}
For $F\in1+x\Z[[x]]$ there are integers $c_d$ such that
$F=\prod_{d\geq1}(1-x^d)^{-c_d}$.
To construct them, remove the coefficient of degree $d$ by multiplying
the residual series by $(1-x^d)^{c_d}$, proceeding in increasing degree.
All these operations preserve integer coefficients, including when $c_d<0$.
Taking the formal logarithmic derivative of $F=F_L$ yields
\[
 L(n)=[x^n]\frac{xF'(x)}{F(x)}=\sum_{d\mid n}d c_d.
\]
Thus $L(n)\in\Z$.
If $m=p^sb$ with $p\nmid b$, the divisors of $pm$ not dividing $m$
are precisely $p^{s+1}e$ with $e\mid b$. Consequently
\[
 L(pm)-L(m)=p^{s+1}\sum_{e\mid b} e c_{p^{s+1}e},
\]
which is divisible by $p$. Apply Theorem~\ref{weighted} with $W(n)=1$.
\end{proof}

\begin{corollary}[Two integral inputs]\label{twoinputs}
Let $L,W:\Z_{>0}\to\Q$, with $F_W$ defined as in \eqref{input}.
If both $F_L$ and $F_W$ have integer coefficients,
then \eqref{weighted-equation} has a unique integral solution.
\end{corollary}
\begin{proof}
The Euler-product argument just given applies separately to $L$ and $W$.
It makes both sequences integer-valued and gives
$W(pm)-W(m)\in p^{v_p(m)+1}\Z$, as well as the needed congruence for $L$.
Theorem~\ref{weighted} applies.
\end{proof}

\begin{corollary}[Divisor sums and geometric dilations]\label{instances}
The defining series of A163921, A157675 and A159593 have integer coefficients.
More generally, $L(n)=a^n$ works for every integer $a$.
For $a=0$ the solution is $(1-x)^{-1}$; for $a=1$ it is the Catalan series.
\end{corollary}
\begin{proof}
For $L(n)=\sigma(n)$, take $c_d=1$ for every $d$; then
$F_L=\prod_{d\geq1}(1-x^d)^{-1}\in\Z[[x]]$.
For every integer $a$, the choice $L(n)=a^n$ has
$F_L=(1-ax)^{-1}$. Thus A157675 ($a=2$) and A159593 ($a=3$)
have integer coefficients~\cite{oeis157675,oeis159593}.
For $a=0$ the solution is $(1-x)^{-1}$; for $a=1$ it satisfies
$G=(1-xG)^{-1}$, equivalently $G=1+xG^2$, the Catalan equation.
\end{proof}

\begin{corollary}[An unasserted converse fails]\label{converse}
For every prime $q$, set $L_q(n)=q$ if $q^2\mid n$ and $L_q(n)=0$
otherwise. Its Hanna fixed point is integral, but
\[
 F_{L_q}(x)=(1-x^{q^2})^{-1/q},\qquad [x^{q^2}]F_{L_q}=1/q.
\]
Thus integrality of the fixed point does not imply integrality of the input.
This converse is not asserted in the posing source.
\end{corollary}
\begin{proof}
For a prime $p\ne q$, divisibility by $q^2$ is unchanged on replacing $m$
by $pm$, so $L_q(pm)=L_q(m)$. For $p=q$ both values are divisible by $q$.
Corollary~\ref{original} applies under its weaker hypothesis. The displayed
input follows by summing $\sum_{k\geq1}q x^{q^2k}/(q^2k)$.
\end{proof}

\begin{corollary}\label{shifted}
For every integer $a$, the equation
\begin{equation}\label{shifted-equation}
 A(x)=1+x\exp\left(\sum_{n\geq1}\frac{(A(a^nx)-1)^n}{n}\right)
\end{equation}
has a unique solution in $\Q[[x]]$, and it has integer coefficients.
This proves the conjectures of A156907 ($a=2$) and A156908 ($a=3$)~\cite{oeis156907}.
\end{corollary}
\begin{proof}
Any solution has the form $A=1+xG$ with $G(0)=1$.
Since $(A(a^nx)-1)^n=a^{n^2}x^nG(a^nx)^n$, the equation for $G$
is \eqref{weighted-equation} with $L(n)=a^n$ and $W(n)=a^{n^2}$.
The first congruence in \eqref{congruences} follows from Fermat's theorem.
For the second write $s=v_p(m)$.
If $p\mid a$ and $a\ne0$, both $a^{m^2}$ and $a^{p^2m^2}$ are divisible
by $p^{s+1}$, since $m^2\geq s+1$; the case $a=0$ is immediate.
If $p\nmid a$, Euler's theorem modulo $p^{s+1}$ applies because
$m^2(p^2-1)$ is divisible by $p^s(p-1)$.
Thus $a^{p^2m^2}\equiv a^{m^2}\pmod{p^{s+1}}$, as required.
Theorem~\ref{weighted} proves the assertion, including $a\leq0$.
\end{proof}

\begin{corollary}[Sparse substitution]\label{sparse}
For any integer $a$, the equation
\begin{equation}\label{sparse-equation}
 A(x)=\exp\left(\sum_{n\geq1}\frac{x^n}{n}A(a^nx^n)\right)
\end{equation}
has a unique integral solution, including A179471 when $a=2$~\cite{oeis179471}.
Its logarithm starts with $xA(ax)+x^2A(a^2x^2)/2$.
\end{corollary}
\begin{proof}
For integral $B$ with $B(0)=1$, put $K_B=\sum_{n\geq1}x^nB(a^nx^n)/n$.
The second sum in the analogue of \eqref{pairing} has brackets
\[
 B(a^{pm}x^{pm})-B(a^m x^{pm}).
\]
For $s=v_p(m)$ one has $a^{pm}\equiv a^m\pmod{p^{s+1}}$:
if $p\mid a$, use $m\geq s+1$ and handle $a=0$ directly;
otherwise use Euler's theorem, since $m(p-1)$ is divisible by $p^s(p-1)$.
Every coefficient of the bracket is therefore divisible by $p^{s+1}$.
Lemma~\ref{dwork} proves that $\exp(K_B)$ is integral.
If $B-C$ is divisible by $x^d$, the $n$th summand of $K_B-K_C$
is divisible by $x^{n(d+1)}$, so the same fixed-point argument proves
existence and uniqueness. Taking the formal logarithm proves the final claim.
\end{proof}

An Euler-product proof uses Manyama's identity~\cite{manyama}.
Write $A(x)=\sum_{j\geq0}b_jx^j$. Then
\[
 \exp\left(\sum_{n\geq1}\frac{x^n}{n}A(a^nx^n)\right)
 =\prod_{j\geq0}(1-a^jx^{j+1})^{-b_j}.
\]
Indeed, expand $A$ and sum
$\sum_{n\geq1}(a^jx^{j+1})^n/n=-\log(1-a^jx^{j+1})$.
The coefficient of degree $N$ in this product depends only on
$b_0,\ldots,b_{N-1}$. Starting with $b_0=1$, integral binomial
expansions therefore prove integrality and uniqueness recursively,
for every integer $a$. The identity supplies the Euler product; its
triangular coefficient dependence supplies the fixed-point conclusion.

The example line of A179471 prints
``$\log(A(x))=A(2x)+\cdots$''. Its first term needs the factor $x$:
the printed right-hand side has constant coefficient $1$, whereas
$\log A$ has constant coefficient $0$. The corrected expansion is
\[
 \log A(x)=xA(2x)+\frac{x^2}{2}A(4x^2)+\frac{x^3}{3}A(8x^3)+\cdots.
\]
The defining equation and the PARI expression in that entry already agree
with this correction. The numerical asymptotic constant listed there is
not covered by the present argument.

\begin{corollary}[Nonnegative coefficients]\label{nonnegative}
If $L(n),W(n)\geq0$, the unique rational solution of
\eqref{weighted-equation} has nonnegative coefficients.
All six OEIS series covered here have nonnegative integer coefficients.
\end{corollary}
\begin{proof}
Dilation by a nonnegative number, positive integer powers, addition,
and formal exponentiation preserve nonnegative coefficients.
Thus every iterate starting at $1$ is nonnegative, and so is its
coefficientwise limit. The same argument applies to the sparse equation
when $a\geq0$. The shifted series has $A=1+xG$ with
$L(n)=a^n\geq0$ and $W(n)=a^{n^2}\geq0$ for its positive source bases.
Integrality was proved above.
\end{proof}

\section{Coverage and verification}

The posing entries give the following coverage. Their definitions, comments,
and comparisons are separated from the asymptotic assertion.
\begin{center}
\small
\begin{tabular}{@{}>{\raggedright\arraybackslash}p{0.19\textwidth}>{\raggedright\arraybackslash}p{0.40\textwidth}>{\raggedright\arraybackslash}p{0.31\textwidth}@{}}
\toprule
Source & Statement & Coverage\\
\midrule
A163921, comment & Conditional integrality for rational $L$ & Reproved; prior proof via Li--Zhong--Zhou (2026)\\
A163921, definition & $L(n)=\sigma(n)$ & Reproved; same reduction\\
A157675 & $L(n)=2^n$ & Reproved; same reduction\\
A157675, A159593, comparison & Catalan equation & Prior identity; derived in Corollary~\ref{instances}\\
A159593, definition & $L(n)=3^n$ & Reproved; same reduction\\
A156907, A156908 & Shifted integrality conjectures & Reproved; prior proof via Li--Zhong--Zhou (2026)\\
A179471, definition & Sparse substitution, $a=2$ & Proved; Euler identity of Manyama (2017)\\
A179471, example & Logarithm without initial $x$ & Refuted-and-corrected, expansion above\\
A179471, formula & Asymptotic constant & Not covered\\
All six entries, keyword & Nonnegative coefficients & Proved, Corollary~\ref{nonnegative}\\
Not posed & Two integral inputs & Proved; also follows by the same reduction\\
Not asserted in source & Converse of A163921 & Fails for $L_2$, Corollary~\ref{converse}\\
\bottomrule
\end{tabular}
\end{center}

``Prior proof via'' denotes the explicit deduction from the polynomial
theorem above; Li--Zhong--Zhou do not name these six target entries.
Their explicit implicit-equation application is A163138.

In the pinned
\href{https://github.com/the-omega-institute/trureturing/blob/91f44a2b8e396fd6374dd8ae047c158afbbebba4/D5/S1/Recurrence/Algebraic/HannaSelfSubstitutionIntegrality.lean}{main Lean file},
\texttt{result} in the namespace
\texttt{HannaSelfSubstitutionIntegrality}
proves existence and integrality of every solution under the integral-input
hypothesis of \eqref{input}, for every rational sequence $L$.
Its positive-index sum \texttt{logSum} has degree-$N$ coefficient
$\sum_{n=1}^N[x^{N-n}]f(n)/n$, so the index-zero value of $L$ is unused.
The definition \texttt{SolvesHanna} expresses precisely that fixed-point equation;
\texttt{IntegralSeries} means that every rational coefficient is an integer.
In the pinned
\href{https://github.com/the-omega-institute/trureturing/blob/91f44a2b8e396fd6374dd8ae047c158afbbebba4/D5/S1/Recurrence/Algebraic/HannaSelfSubstitutionIntegralityA157675.lean}{specialization Lean file},
\texttt{result} in the namespace
\texttt{HannaSelfSubstitutionIntegralityA157675}
proves the $L(n)=2^n$ case. The same main file's \texttt{geometric\_input}
proves integrality of $F_L$ for every integer base $a$ and $L(n)=a^n$;
it is not a theorem about all those fixed points by itself.
The full namespaces begin with \texttt{D5.S1.Recurrence.Algebraic.}
Uniqueness, the weighted and two-input theorems, the shifted and sparse
extensions, and the converse obstruction have the written proofs given here.

The accompanying \texttt{verify.py} compares rational coefficient recursion with
formal iteration and checks the available OEIS coefficient lists. Its finite
scope and exact coefficients are recorded in \texttt{verification.json}.
These coefficient checks are distinct from the all-degree proofs above.

\begin{thebibliography}{99}
\bibitem{lzz} J. Li, Y. Zhong and Y. Zhou,
\emph{Integral Exponential Series: Fibonacci Weights and Implicit Equations},
arXiv:2611.16283v1 (12 November 2026), Theorem 1.2,
equations (1.5)--(1.6), p.~2, and Section 5, pp.~7--8.
\url{https://arxiv.org/abs/2611.16283v1}.
\bibitem{bala} P. Bala, integrality proof for integer square-exponential
weights, comment in \emph{OEIS A155200}, 25 December 2019.
\url{https://oeis.org/A155200}.
\bibitem{jmw} S. Jaidee, P. Moss and T. Ward,
\emph{Time-changes preserving zeta functions}, arXiv:1809.09199,
Theorem 2 and Example 7.
\url{https://arxiv.org/abs/1809.09199}.
\bibitem{manyama} S. Manyama, generalized Euler-transform identity,
FORMULA block in \emph{OEIS A266964}, 14 November 2017.
\url{https://oeis.org/A266964}.
\bibitem{oeis163921} P. D. Hanna, \emph{OEIS A163921}, unpaginated comment and
definition, contributed 1 September 2009, revision \#2, 30 March 2012.
\url{https://oeis.org/A163921}.
\bibitem{oeis157675} P. D. Hanna, \emph{OEIS A157675}, comment and Catalan
comparison, contributed 6 March 2009, revision \#2, 30 March 2012.
\url{https://oeis.org/A157675}.
\bibitem{km} C. Krattenthaler and T. W. M\"uller,
\emph{Truncated versions of Dwork's lemma for exponentials of power series
and $p$-divisibility of arithmetic functions}, Adv. Math. \textbf{283} (2015),
489--529. \url{https://doi.org/10.1016/j.aim.2015.05.018}.
The cited version is arXiv:1412.7014v2, Lemma 1, p.~3.
\url{https://arxiv.org/abs/1412.7014v2}.
\bibitem{ps} J. Pomerat and A. Straub,
\emph{Criteria for the integrality of $n$th roots of power series},
arXiv:2406.12010v1 (17 June 2024), Theorem 1.1 and Corollary 4.2.
\url{https://arxiv.org/abs/2406.12010v1}.
\bibitem{bgw} J. Byszewski, G. Graff and T. Ward,
\emph{Dold sequences, periodic points, and dynamics},
Bull. London Math. Soc. \textbf{53} (2021), 1263--1298,
Theorem 2.5 and Lemma 2.6. \url{https://doi.org/10.1112/blms.12531}.
\bibitem{oeis159593} P. D. Hanna, \emph{OEIS A159593}, definition,
contributed 2 May 2009, revision \#2, 30 March 2012.
\url{https://oeis.org/A159593}.
\bibitem{oeis156907} P. D. Hanna, \emph{OEIS A156907 and A156908},
integrality comments, contributed 4 March 2009, revision \#2, 30 March 2012.
\url{https://oeis.org/A156907}; \url{https://oeis.org/A156908}.
\bibitem{oeis179471} P. D. Hanna, \emph{OEIS A179471}, definition and formula,
contributed 15 July 2010, revision \#4, 4 October 2013.
\url{https://oeis.org/A179471}.
\end{thebibliography}
\noindent\small\textbf{License.} Creative Commons Attribution 4.0 International
(CC BY 4.0). \url{https://creativecommons.org/licenses/by/4.0/}.
\end{document}
